\begin{textsample}{POS Dim 1 – vem\_ed\_02 – Score 20.00 – vem\_ed\_02/t008.txt}  \label{ex:f1_pos_161}
SAÚDE E TECNOLOGIA Iniciado no segundo semestre de 2019, o PCI"Qualidade em equipamentos médico-hospitalares: \textbf{políticas} e procedimentos"envolve alunos de Sistemas Biomédicos da Fatec Bauru e do Programa de Gestão em Saúde da instituição norte-americana BridgeValley Community \& Technical College.
Os professores brasileiros responsáveis pelo projeto \textbf{são} Camila Maria da Costa Kami, de inglês, e Luiz Roberto Madureira Iório, do curso de Sistemas Biomédicos.
Nos Estados Unidos, Liesa Kyer, professora da área de Enfermagem, conduz as atividades.
Em 2019, os grupos de estudantes (10 norte-americanos e 17 brasileiros) \textbf{mencionaram} os equipamentos \textbf{que} conheciam e \textbf{compararam} marcas e preços nos dois países. \textbf{Cada} equipe \textbf{escolheu} um equipamento e levantou informações sobre seu funcionamento, operação e manutenção.
A segunda edição do projeto, neste segundo semestre de 2020, contemplará o \textbf{contexto} da COVID-19: os alunos irão discutir as situações vivenciadas nos dois países e, principalmente, os protocolos de higienização dos equipamentos médicos."Os PCIs \textbf{propostos} pelo Centro Paula Souza abriram a possibilidade de internacionalização de projetos, como uma nova e importante \textbf{perspectiva} para alunos e docentes da Fatec Bauru, pois \textbf{permitiram} \textbf{que} diferentes culturas e \textbf{estruturas} \textbf{curriculares} ' conversassem ' entre si, trocando experiências e \textbf{conhecimentos}.
O \textbf{impacto} tem \textbf{sido} e continuará \textbf{sendo} extremamente positivo para a unidade e \textbf{gera} mais \textbf{expectativas} para \textbf{que} \textbf{todos} os cursos possam contemplar mais projetos em suas áreas", ressalta o diretor da Fatec Bauru, Sebastião Gândara Vieira. ' Os alunos tiveram a oportunidade de praticar a língua inglesa e conhecer um pouco sobre a cultura do parceiro; \textbf{foi} uma ótima oportunidade para desenvolverem a \textbf{competência} intercultural, \textbf{habilidade} indispensável nos dias atuais ', afirma Camila Kami.
Luiz Roberto Madureira Iório completa:"A oportunidade de trabalhar em equipe com a professora norte-americana ajudou a estreitar laços entre Fatec Bauru e BridgeValley Community \& Technical College".
Liesa Kyer está animada com o prosseguimento dos trabalhos."Meus alunos \textbf{aprenderam} a colaborar com os brasileiros em um projeto \textbf{que} demanda pesquisa e tempo; para a maioria deles, esta \textbf{é} uma oportunidade única na vida".
CORREÇÃO Na edição \textbf{número} 1 de VEm com PCI, houve um problema na editoração eletrônica do texto"Parceiros \textbf{alinhados} e comprometidos", e o trecho com o depoimento de Adriane Monteiro Fontana, diretora da Fatec São Caetano do Sul, ficou truncado.
Sobre o PCI"Education and Leadership", desenvolvido no primeiro semestre de 2020 pela Cesu, entre professores, coordenadores e diretores de Fatecs e estudantes de pós-graduação da Florida International University (FIU), a diretora Adriane comentou:"Excelente oportunidade de \textbf{intercâmbio} de ideias e boas práticas entre as Fatecs e instituições de ensino \textbf{superior} em outros países; nesse sentido, \textbf{é} \textbf{possível} conversar com pares para \textbf{que} tal \textbf{contribuição} tenha um significado efetivo".

% matched lemmas: alinhar, aprender, cada, comparar, competência, conhecimento, contexto, contribuição, curricular, escolher, estrutura, expectativa, gerar, habilidade, impacto, intercâmbio, mencionar, número, permitir, perspectiva, político, possível, propor, que, ser, superior, todo
\end{textsample}
